\documentclass[11pt,letter]{article}
\usepackage[top=1.00in, bottom=1.0in, left=1.1in, right=1.1in]{geometry}
\renewcommand{\baselinestretch}{1.5}
\usepackage{graphicx}
\usepackage{natbib}
\usepackage{amsmath}
\usepackage{hyperref}

\def\labelitemi{--}
\begin{document}

\title{Core mounting protocol} 

\setlength{\parindent}{0pt}
\setlength{\parskip}{3pt}
\begin{enumerate}
\item Copy paste the folder you are going to work on to the laptop
\item Navigate into one stand/species folder
\item Create a new folder in this folder, name this folder: ``MeasurementStandSpecies'' (replace the Stand and Species with the ones in the folder name)
\item Species name: CANO, TSME, TSHE, THPL, PSME, ABAM; Stand ID: AV06, AB08, AO03, TO04, AE10, AG05, TA01, PARA, SUNR, AM16
\item Open CooRecorder, a window named `Setup project paths' would pop up, if you don't see it, go to File on the navigation bar, and then click on the Setup project paths.
\item For Project name, type in MORA2024StandSpecies (replace the Stand and Species with the folder name that you are going to work on)
\item For Image files directory tree, click on it and navigate to the folder you are going to work on
\item For .Pos files directory tree, click on it and navigate to the MeasurementStandSpecies folder you created
\item Click Ok \& Close
\item Go to File on the navigation bar again, click on Open image file for new coordinates
\item Once the image opens, there will be a window pop up
\item In the select measurement resolution section, make sure `Measure in milimetres' is selected, then change the new DPI value to 6450.
\item In the Data modes for dendrochronology section, make sure `Measure ring widths or generate Blue data or XRay density data. (which is usually the default)'
\item Click OK.
\item Now you can start adding points to the ring, if you don't remember how to do it, please refer to the tutorial video I sent to you.
\item Tip: Some rings might be wavy, use the multiple points to move around to find the most paralleled section
\item Tip: Some time when you try to open a file it might tell you unable to open the file, then just quit the application and open again, if it still doesn't work, leave a note
\item IMPORTANT: Always start from the very end of the core boundary (see the attached screenshots for some examples)
\item Only do the measurement when you are confident with it, if you are not sure or the scanning are unclear, you can either: (1) Go to the lab and refer to the physical cores; (2)Just leave a note and I will deal with them later
\item Once the measurement is done, go to File, click on Save As, name the .Pos file according to the folder name of this specific sample (E.g., PSME\_AB08\_4125(1))
\item Fill out the excel sheet whenever one sample is done, record: how many rings measured, and notes need to pay attention later.
	\end{enumerate}
\end{document}